\documentclass[10pt,a4paper]{article}
\usepackage[T1]{fontenc}
\usepackage{lmodern,amsmath,amssymb,amsthm,microtype,xcolor,geometry,hyperref,enumitem}
\geometry{margin=19mm}
\definecolor{softgray}{gray}{0.94}
\hypersetup{hidelinks,pdftitle={SGA 2 Expose VIII Corollary 2.2 statement},pdfauthor={Interlanguage project}}
\newcounter{sourcecorollary}[section]
\renewcommand{\thesourcecorollary}{\thesection.\arabic{sourcecorollary}}
\newcommand{\cF}{\mathcal F}
\newcommand{\markedfootnote}[2]{%
  \begingroup\renewcommand{\thefootnote}{#1}\footnote{#2}\endgroup}
\newenvironment{sourcecorollary}[1]{%
  \refstepcounter{sourcecorollary}%
  \par\medskip\noindent
  \textbf{Corollary~\thesourcecorollary\markedfootnote{(3)}{#1}.}\enspace
  \itshape
}{\par\medskip}
\begin{document}
\begin{center}
  {\Large\bfseries SGA 2: Local Cohomology and Lefschetz Theorems}\par\smallskip
  {\large Expos\'e VIII: The Finiteness Theorem}\par
  \small Bounded source-aligned working unit --- 19 July 2026
\end{center}
\noindent\colorbox{softgray}{\parbox{0.95\linewidth}{\footnotesize\textbf{Authority.}
Corrected French lines 2661--2668; original printed p.~89; physical
source-PDF p.~79; recomposed running p.~71. The compiled reader is the output
of the same corrected French edition, not an independent original typescript
witness. Line 2669 is blank; Corollary~2.3 begins at line 2670 and is excluded.}}

\setcounter{section}{2}
\setcounter{sourcecorollary}{1}
\begin{sourcecorollary}{\textit{Editor's note.} Strictly speaking, this is a
corollary of the proof that follows, not of the statement. The implication
\textup{c)}$\Rightarrow$\textup{a)} is tautological. The converse is not, but
follows from the proof. More precisely, as below, one covers $X$ by open
subsets embeddable in regular schemes; as explained below, this permits one to
reduce to the case where $X=\operatorname{Spec}(A)$ is affine and regular and
$\cF=\widetilde M$, where $M$ is an $A$-module of finite projective dimension.
In this case it is shown that conditions \textup{a)} and \textup{c)} are
equivalent to the dual conditions \textup{a$'$)} and \textup{c$'$)}. One then
shows that \textup{c$'$)} implies condition \textup{d)} (see below), which in
turn implies \textup{a$'$)}. See the discussion following~2.4.}\label{VIII.2.2}
Under the hypotheses of the preceding theorem, condition \textup{a)} is
equivalent to:
\begin{enumerate}[label=\textup{\alph*)},start=3,leftmargin=2.2em]
\item for every $x\in U$ such that $c(x)=1$, one has
  $H^{i-1}(\cF_x)=0$.
\end{enumerate}
\end{sourcecorollary}
\end{document}
